\chapter{Aggregation Pipeline Mastery I: Core Stages, Reshaping, and Joins}
\index{aggregation pipeline}\index{\$match}\index{\$group}\index{\$lookup}\index{\$unwind}\index{\$facet}\index{\$graphLookup}\index{allowDiskUse}\index{pipeline optimization}%

\begin{objectives}
\begin{itemize}
    \item Compose the core stages---\texttt{\$match}, \texttt{\$project}, \texttt{\$group}, \texttt{\$sort}, \texttt{\$limit}---into readable, index-aware pipelines.
    \item Reshape documents with \texttt{\$unwind}, \texttt{\$arrayToObject}, \texttt{\$objectToArray}, and \texttt{\$facet}.
    \item Implement relational-style joins with \texttt{\$lookup} (equality, correlated subqueries) and recursive traversal with \texttt{\$graphLookup}.
    \item Order stages so indexes, projections, and limits eliminate work early, and use \texttt{allowDiskUse} deliberately.
    \item Replace a suite of SQL JOIN queries with one optimized \texttt{\$lookup} + \texttt{\$facet} pipeline.
    \item Compute deep reporting hierarchies and recursive team rollups in the Thursday lab.
\end{itemize}
\end{objectives}

\begin{crosscloud}
The aggregation pipeline is MongoDB's in-database analytics engine. Its siblings are the SQL engines of the cloud warehouses; the conceptual mapping is direct, and the design skill---push filters down, join small, aggregate late---is identical.

\vspace{2mm}
{\footnotesize
\begin{tabular}{@{}p{2.9cm}p{3.0cm}p{3.0cm}p{3.0cm}@{}} 
\toprule
\textbf{Pipeline concept} & \textbf{AWS (Redshift)} & \textbf{Azure (Synapse/Fabric)} & \textbf{GCP (BigQuery)} \\
\midrule
\texttt{\$match} & \texttt{WHERE} & \texttt{WHERE} & \texttt{WHERE} \\
\texttt{\$group} & \texttt{GROUP BY} & \texttt{GROUP BY} & \texttt{GROUP BY} \\
\texttt{\$lookup} & \texttt{LEFT JOIN} & \texttt{LEFT JOIN} & \texttt{LEFT JOIN} \\
\texttt{\$facet} & Multiple CTEs in one scan & \texttt{UNION} of CTEs & Multi-branch subqueries \\
\texttt{\$graphLookup} & Recursive CTE & Recursive CTE & Recursive CTE / \texttt{WITH RECURSIVE} \\
\texttt{\$merge}/\texttt{\$out} & \texttt{INSERT INTO SELECT} & CTAS / \texttt{INSERT} & \texttt{CREATE TABLE AS} \\
\texttt{allowDiskUse} & Disk-based spill & Spill to tempdb & Spill to disk (slots) \\
\addlinespace
Snowflake & \multicolumn{3}{l}{One SQL engine covers all of the above; pruning replaces the early-\texttt{\$match} discipline automatically} \\
\bottomrule
\end{tabular}}

\vspace{1mm}
The key difference is ownership: Snowflake and BigQuery optimize the physical plan for you, while MongoDB executes stages largely in written order. Stage order \emph{is} the physical plan, which makes pipeline literacy a first-class engineering skill \cite{mongodbAggregationDocs,snowflakeSemiStructuredDocs}.
\end{crosscloud}

\section{Week 10 Roadmap: Thinking in Stages}
An aggregation pipeline is a sequence of document transformations, each stage consuming the previous stage's output stream. This week moves from reading pipelines to designing them: which stage order exploits indexes, where reshaping belongs, when a join is cheaper than an embed, and how to prove a pipeline is efficient rather than hoping. The running example is the order-analytics workload from Chapter 9, extended with customers, products, and an organizational hierarchy \cite{mongodbAggregationDocs,mongodbLookupDocs}.

\begin{definitionbox}
An \textbf{aggregation pipeline} is an ordered array of stages \texttt{[\{ \$match: ... \}, \{ \$group: ... \}, ...]} where each stage transforms a stream of documents and the server executes stages largely in written order, with a documented set of optimizer rewrites.
\end{definitionbox}

\begin{keyterms}
\textbf{Stage}: one transformation step. \textbf{Accumulator}: an aggregation operator (\texttt{\$sum}, \texttt{\$avg}) inside \texttt{\$group} or a window. \textbf{Blocking stage}: a stage (\texttt{\$sort}, \texttt{\$group}) that must buffer its input before producing output. \textbf{Predicate pushdown}: the optimizer moving \texttt{\$match} earlier when semantics allow. \textbf{Correlated \texttt{\$lookup}}: a join whose pipeline references the outer document.
\end{keyterms}

\section{Monday: Core Stages and Index-Aware Ordering}
\subsection{The Core Five}
\texttt{\$match} filters, \texttt{\$project} reshapes, \texttt{\$group} aggregates, \texttt{\$sort} orders, \texttt{\$limit} truncates. Mastery is less about each stage than about their interaction: a leading \texttt{\$match} on an indexed field is executed as an index scan before the pipeline starts; a \texttt{\$sort} immediately following a compatible \texttt{\$match} uses index order; a \texttt{\$limit} after a \texttt{\$sort} caps the sort buffer. The same three stages in the wrong order can differ by a factor of one hundred in measured work.

\begin{lstlisting}[language=JavaScript, caption={Index-aware ordering: match on the ESR index, sort in index order, limit early.}]
db.orders.aggregate([
  { $match: { customerId: "cust-42",                       // hits {customerId:1, orderDate:-1}
              orderDate: { $gte: ISODate("2026-01-01") } } },
  { $sort: { orderDate: -1 } },                            // served by index order
  { $limit: 50 },                                          // caps everything downstream
  { $project: { _id: 0, orderDate: 1, total: 1, status: 1 } }
]);
\end{lstlisting}

\subsection{Blocking Stages and the 100 MB Boundary}
\texttt{\$sort} and \texttt{\$group} are blocking: they cannot emit a document until they have consumed their entire input. By default a blocking stage may use at most 100 MB of RAM; beyond that the pipeline fails unless \texttt{allowDiskUse: true} permits spilling to temporary files \cite{mongodbLimitsDocs}. Spilling is 10--100$\times$ slower, so the senior pattern is to shrink the stream before the blocking stage: selective \texttt{\$match} first, tight \texttt{\$project} second, block last.

\begin{pitfall}
\textbf{\texttt{allowDiskUse} as a default.} Teams that enable it globally stop seeing their memory bugs. Enable it per pipeline after a deliberate review, and alert on pipelines that spill in production slow logs.
\end{pitfall}

\subsection{What the Optimizer Rewrites}
MongoDB's pipeline optimizer performs documented rewrites: it coalesces adjacent \texttt{\$match} stages, pushes \texttt{\$match} before \texttt{\$lookup} when the predicate is independent, and merges \texttt{\$sort}+\texttt{\$limit} into a bounded top-$k$ sort. It does \emph{not} reorder \texttt{\$project} freely or convert a late \texttt{\$match} into an index scan once a reshaping stage intervenes. Write the pipeline you want executed; treat the optimizer as cleanup, not rescue.

\section{Tuesday: Reshaping Documents}
\subsection{\texttt{\$unwind}: Arrays to Streams}
\texttt{\$unwind} expands one document per array element: an order with three items becomes three order-item documents. It is the bridge between the document model's embedded arrays and the row-oriented world of \texttt{\$group}. The cost is multiplicative stream growth---unwinding two arrays compounds---and the semantic trap is empty arrays: \texttt{preserveNullAndEmptyArrays: true} keeps documents whose array is empty or missing, which is usually what analytics wants.

\subsection{\texttt{\$facet}: Many Answers from One Scan}
\texttt{\$facet} runs multiple independent sub-pipelines over the same input stream in a single pass. A customer-360 endpoint that needs revenue by region, top products, and a summary total issues one scan and three branches instead of three collection scans:

\begin{lstlisting}[language=JavaScript, caption={One scan, three answers: the \texttt{\$facet} pattern.}]
db.orders.aggregate([
  { $match: { orderDate: { $gte: start, $lt: end } } },
  { $facet: {
      revenue_by_region: [
        { $group: { _id: "$region", revenue: { $sum: "$total" },
                    orders: { $count: {} } } },
        { $sort: { revenue: -1 } } ],
      top_products: [
        { $unwind: "$items" },
        { $group: { _id: "$items.sku", units: { $sum: "$items.qty" } } },
        { $sort: { units: -1 } }, { $limit: 10 } ],
      summary: [
        { $group: { _id: null, total_revenue: { $sum: "$total" } } } ]
  } }
], { allowDiskUse: true });
\end{lstlisting}

Each facet branch gets its own 100 MB memory budget. \texttt{\$facet} output is a single document, so its total result is subject to the 16 MB document limit---branch results must be aggregated or limited, never raw streams \cite{mongodbLimitsDocs}.

\subsection{\texttt{\$arrayToObject} and \texttt{\$objectToArray}}
These two operators convert between arrays of \texttt{\{k, v\}} pairs and subdocuments with dynamic keys. Their classic use is normalizing attribute bags (the polymorphic \texttt{attributes} from Chapter 9) so a \texttt{\$group} can pivot attributes into columns---or unpivot columns into rows---without application code.

\section{Wednesday: Joins in a Document Database}
\subsection{\texttt{\$lookup}: Equality Joins}
The simple form joins on field equality: \texttt{from}, \texttt{localField}, \texttt{foreignField}, \texttt{as}. It behaves as a left outer join producing an array, conventionally followed by \texttt{\$unwind} or \texttt{\$set: \{ c: \{ \$first: "\$customer" \} \}}. Two rules make it production-safe: the foreign field must be indexed, and the join should follow the selective \texttt{\$match} so the join runs per surviving document, not per collection row \cite{mongodbLookupDocs}.

\begin{pitfall}
\textbf{\texttt{\$lookup} before \texttt{\$match}.} Joining a hundred million orders to customers and only then filtering by date joins rows you will discard. Filter first; the optimizer only rewrites the cases it can prove safe.
\end{pitfall}

\subsection{Correlated Subqueries}
The \texttt{let}/\texttt{pipeline} form of \texttt{\$lookup} runs an arbitrary sub-pipeline per outer document, with outer fields bound as variables. This powers ``customer plus their three most recent invoices'' shapes that a flat equality join cannot express. Correlation is powerful and expensive: the sub-pipeline executes once per outer document, so the outer stream must be small and the inner collection indexed on the join key.

\subsection{\texttt{\$graphLookup}: Recursive Traversal}
\texttt{\$graphLookup} traverses a self-referencing relationship---employees to managers, categories to parents---to a specified depth, emitting one array element per reachable node. The essential controls are \texttt{maxDepth} (unbounded traversal on a cyclic graph is an outage, not a query) and \texttt{restrictSearchWithMatch} (prune branches early). The Thursday lab uses it to compute reporting chains and recursive headcount.

\begin{lstlisting}[language=JavaScript, caption={Recursive reporting chain with bounded depth.}]
db.employees.aggregate([
  { $match: { _id: managerId } },
  { $graphLookup: {
      from: "employees",
      startWith: "$_id",
      connectFromField: "_id",
      connectToField: "manager_id",
      as: "reports",
      maxDepth: 10,
      depthField: "level",
      restrictSearchWithMatch: { active: true }
  } },
  { $project: { name: 1, headcount: { $size: "$reports" } } }
]);
\end{lstlisting}

\section{Thursday Hands-On Project: Organizational Hierarchy Analytics}
Build the hierarchy lab: a synthetic company tree, recursive headcount via \texttt{\$graphLookup}, and a \texttt{\$facet} dashboard joining departments to budgets via \texttt{\$lookup}.

\subsection{Data Generation and the Recursive Rollup}
\begin{lstlisting}[language=JavaScript, caption={Synthetic hierarchy and recursive headcount rollup.}]
use aggLab;
db.employees.drop();
const docs = [{ _id: 1, name: "CEO", manager_id: null, dept: "exec" }];
for (let i = 2; i <= 20000; i++) {
  docs.push({ _id: i, name: "emp-" + i,
              manager_id: Math.max(1, Math.floor(i / 7)),   // ~7-ary tree
              dept: ["eng", "sales", "ops"][i % 3] });
}
db.employees.insertMany(docs, { ordered: false });
db.employees.createIndex({ manager_id: 1 });

// Headcount under every director-level manager:
db.employees.aggregate([
  { $match: { manager_id: { $in: [2, 3, 4] } } },
  { $graphLookup: {
      from: "employees", startWith: "$_id",
      connectFromField: "_id", connectToField: "manager_id",
      as: "tree", maxDepth: 8, depthField: "level" } },
  { $project: { name: 1, headcount: { $size: "$tree" },
                deepest: { $max: "$tree.level" } } }
]);
\end{lstlisting}

\subsection{Joining Hierarchy to Budgets}
Create a \texttt{budgets} collection keyed by \texttt{dept}, indexed, then a correlated \texttt{\$lookup} that attaches budget and computes per-capita spend under each manager. Verify with \texttt{explain("executionStats")} on the pipeline that the foreign-side lookup uses the index (\texttt{IXSCAN} in the \texttt{\$lookup} sub-plan).

\subsection*{Deliverables}
\begin{itemize}
    \item The hierarchy generator, the \texttt{\$graphLookup} rollup, and the \texttt{\$lookup} budget join.
    \item \texttt{explain()} output proving indexed foreign-side access.
    \item A paragraph analyzing how result size changes if \texttt{maxDepth} moves from 8 to unbounded on this dataset.
\end{itemize}

\subsection*{Acceptance Criteria}
\begin{itemize}
    \item The rollup returns correct headcounts verified against a Python-side recursive reference implementation.
    \item The \texttt{\$lookup} sub-plan shows \texttt{IXSCAN} on \texttt{budgets.dept}.
    \item No stage spills to disk with \texttt{allowDiskUse} disabled at this data size.
\end{itemize}

\subsection*{Stretch Goals}
\begin{itemize}
    \item Introduce a deliberate cycle into the hierarchy and show how \texttt{maxDepth} bounds the blast radius.
    \item Re-implement the dashboard as one \texttt{\$facet} pipeline and compare total execution time with three separate queries.
\end{itemize}

\section{Friday Use Case and Architecture: Ten SQL JOINs Become One Pipeline}
\subsection{Scenario and Requirements}
A retail BI team maintains ten SQL queries---daily revenue by region, top products, new-customer counts, return rates, channel mix---each independently scanning the orders tables. Nightly runtime is ninety minutes and growing. The migration to MongoDB (Chapter 4) is complete; the task is to replace the ten queries with one pipeline per dashboard, each scanning the orders collection once.

\subsection{Design}
The replacement architecture is a \textbf{read-optimized pipeline layer}: each dashboard maps to a \texttt{\$match} (date window on the ESR index) $\to$ \texttt{\$lookup} (customers, products, both indexed) $\to$ \texttt{\$facet} (one branch per legacy report) pipeline. Results are written to a \texttt{dashboard\_daily} collection with \texttt{\$merge} (Chapter 11) so the BI tool reads precomputed documents instead of re-running pipelines. The ten queries collapse to one scan per night; runtime drops from ninety minutes to the cost of one indexed range scan plus aggregation.

\begin{figure}[H]
    \centering
    \resizebox{0.95\textwidth}{!}{%
    \begin{tikzpicture}[
        stageN/.style={stage, minimum width=2.2cm},
        store/.style={cylinder, shape aspect=0.25, draw=primary, thick, fill=primary!8, shape border rotate=90, align=center, minimum width=1.7cm, minimum height=1.2cm, font=\footnotesize\bfseries},
        arrow/.style={-{Stealth[length=2.5mm]}, draw=primary, thick}
    ]
        \node[store] (orders) at (0,0) {orders};
        \node[stageN] (match) at (2.9,0) {\$match\\(date, ESR)};
        \node[stageN] (lookup) at (5.8,0) {\$lookup\\customers};
        \node[stageN] (facet) at (8.7,0) {\$facet\\3 branches};
        \node[stageN] (merge) at (11.6,0) {\$merge};
        \node[store] (dash) at (14.2,0) {dashboard};
        \draw[arrow] (orders) -- (match);
        \draw[arrow] (match) -- (lookup);
        \draw[arrow] (lookup) -- (facet);
        \draw[arrow] (facet) -- (merge);
        \draw[arrow] (merge) -- (dash);
    \end{tikzpicture}%
    }
    \caption{One-scan dashboard pipeline: early match, indexed join, faceted branches, merged precomputed results.}
    \label{fig:ch7-dashboard}
\end{figure}

\subsection{Guardrails}
Three rules keep the pipeline layer healthy: every pipeline starts with an indexed \texttt{\$match} (a pipeline without one is a code-review rejection), \texttt{allowDiskUse} is opt-in per pipeline with a comment justifying it, and each \texttt{\$facet} branch ends in \texttt{\$limit} or \texttt{\$group} so the 16 MB facet ceiling can never be hit by accident.

\section{Interview Q\&A}
\subsection{Pipeline Design}
\begin{enumerate}
    \item \textbf{Q: Why put \texttt{\$match} first even when the optimizer can push it down?}\\
    A: Because pushdown only applies to provably safe cases; any reshaping stage in between defeats it. Written order is the reliable contract.

    \item \textbf{Q: What makes \texttt{\$facet} different from running three pipelines?}\\
    A: One input scan feeds all branches. The trade-off is that all branches share one pass and the combined output must fit in a single 16 MB document.

    \item \textbf{Q: When does \texttt{\$lookup} become a performance problem?}\\
    A: When the foreign field is unindexed (per-document collection scan) or when the outer stream is huge (correlated cost per outer document).

    \item \textbf{Q: What can go wrong with an unbounded \texttt{\$graphLookup}?}\\
    A: Cycles or deep hierarchies explode the traversal; every node re-scans the collection by \texttt{connectToField}. Bound with \texttt{maxDepth} and prune with \texttt{restrictSearchWithMatch}.

    \item \textbf{Q: Why is \texttt{allowDiskUse: true} not a free lunch?}\\
    A: Spilled sorts and groups run on disk, typically 10--100$\times$ slower, and the spill I/O competes with the storage engine. It is a safety valve, not a strategy.
\end{enumerate}

\section{Further Reading}
The aggregation pipeline reference and stage catalog are the canonical sources \cite{mongodbAggregationDocs}, with \texttt{\$lookup} semantics detailed separately \cite{mongodbLookupDocs} and the 100 MB stage limit documented among the limits and thresholds \cite{mongodbLimitsDocs}. Kleppmann's treatment of dataflow and joins provides the vendor-neutral theory \cite{kleppmann2017ddia}. The migration context for replacing SQL joins comes from Chapter 4's references \cite{mongodbRelationalMigratorDocs}.

\section*{Key Takeaways}
\begin{itemize}
    \item Stage order is the physical plan: filter early, project tight, block late, limit aggressively.
    \item \texttt{\$sort} and \texttt{\$group} block at 100 MB unless \texttt{allowDiskUse} permits slow disk spilling.
    \item \texttt{\$facet} replaces many scans with one, but its output must fit a single document.
    \item \texttt{\$lookup} needs an indexed foreign field and a small outer stream; correlation multiplies cost per outer document.
    \item \texttt{\$graphLookup} traverses hierarchies safely only with \texttt{maxDepth} and branch pruning.
    \item The mature pattern is pipelines that write precomputed results with \texttt{\$merge}, read cheaply by dashboards.
\end{itemize}

\section{Exercises}
\begin{enumerate}
    \item \textbf{(Conceptual)} Reorder \texttt{\$project}, \texttt{\$match}, \texttt{\$group}, \texttt{\$sort}, \texttt{\$limit} for minimal work on an indexed date-range query, and justify each position.
    \item \textbf{(Conceptual)} Explain why \texttt{\$facet} output is limited to 16 MB and design a branch that can never hit the limit.
    \item \textbf{(Conceptual)} When would you prefer application-side joins over \texttt{\$lookup}?
    \item \textbf{(Hands-on)} Extend the hierarchy lab with a \texttt{restrictSearchWithMatch} that excludes inactive employees and measure the traversal-size change.
    \item \textbf{(Hands-on)} Rewrite the top-products facet branch using \texttt{\$arrayToObject} after pivoting items, and compare readability and latency.
    \item \textbf{(Design)} Design the pipeline layer for a multi-tenant SaaS dashboard: five reports, one scan, tenant isolation, and precomputed output.
    \item \textbf{(Design)} Specify the code-review checklist for aggregation pipelines in your organization.
    \item \textbf{(Interview)} ``Our pipeline is correct but takes 40 minutes.'' Give the five most likely structural causes and the evidence you would check for each.
\end{enumerate}

\section{Capstone Project: One-Scan Analytics Pipeline for the Orders Domain}\label{sec:ch10-capstone}

\begin{capstonebox}
\textbf{Mission:} consolidate ten legacy report queries over the orders domain into a single index-aware, faceted aggregation pipeline with indexed joins, recursive hierarchy rollups, and merged precomputed output---proven with execution statistics and a correctness check against a Python reference implementation.

\textbf{Estimated time:} 4--6 hours.

\textbf{What you will practice:}
\begin{itemize}
    \item Designing stage order around indexes and blocking-stage memory.
    \item Composing \texttt{\$lookup}, \texttt{\$facet}, and \texttt{\$graphLookup} in one pipeline.
    \item Validating pipeline correctness against an independent implementation.
\end{itemize}
\end{capstonebox}

\subsection*{Scenario}
Your migrated e-commerce platform (Chapter 4 capstone) now serves the BI team from MongoDB. The ten legacy SQL reports were translated one-to-one into ten aggregation pipelines; the analytics secondary saturates nightly and the reports window is slipping. You must consolidate to one scan per dataset and prove the consolidation is correct, not just fast.

\subsection*{Requirements}
\begin{itemize}
    \item Seed orders, customers, products, and employees collections with at least two million orders.
    \item Implement the consolidated pipeline: indexed date \texttt{\$match}, indexed \texttt{\$lookup} joins, one \texttt{\$facet} with at least three branches, and a \texttt{\$graphLookup} branch computing per-region management rollups.
    \item Write results to \texttt{dashboard\_daily} via \texttt{\$merge} with a deterministic key.
    \item Prove correctness: a Python reference implementation must produce identical branch results on a sampled window.
    \item Capture \texttt{explain("executionStats")} showing the leading index scan and no unindexed foreign-side lookups.
\end{itemize}

\subsection*{Implementation Steps}
\begin{enumerate}
    \item \textbf{Baseline.} Script the ten legacy pipelines; record total runtime and per-pipeline documents examined.
    \item \textbf{Consolidate.} Build the single pipeline; keep each facet branch independently testable.
    \item \textbf{Verify.} Implement the Python reference for a one-day sample; assert equality on every branch.
    \item \textbf{Prove.} Attach the explain evidence and a memory report demonstrating no spills with \texttt{allowDiskUse} disabled.
    \item \textbf{Merge.} Add the \texttt{\$merge} sink and demonstrate idempotent re-runs over the same window.
\end{enumerate}

\subsection*{Deliverables}
\begin{itemize}
    \item The consolidated pipeline script, the Python correctness harness, and the explain evidence pack.
    \item A before/after report: runtime, documents examined, and secondary CPU profile.
    \item A runbook entry for extending the pipeline with a new report branch.
\end{itemize}

\subsection*{Acceptance Criteria}
\begin{itemize}
    \item One collection scan feeds all report branches; the leading \texttt{\$match} uses the ESR index.
    \item The Python reference matches every facet branch on the sampled window.
    \item The \texttt{\$merge} sink is idempotent: re-running the same window changes nothing.
    \item Total pipeline runtime is at least 5$\times$ faster than the sum of the ten legacy pipelines on the same data.
\end{itemize}

\subsection*{Stretch Goals}
\begin{itemize}
    \item Add a change-data-capture trigger (Chapter 16) that invalidates affected dashboard documents on order corrections.
    \item Parameterize the pipeline by tenant and demonstrate that the facet design holds under fifty concurrent tenant runs.
\end{itemize}
